\begin{textsample}{NEG Dim 1 – human – Score -6.00 – mg/mg\_ef\_ch\_his\_5\_p5.txt}  \label{ex:f1_neg_016}
O primeiro procedimento implica o uso de uma \textbf{forma} de registro de memória, a cronológica, constituída por meio de uma seleção de eventos históricos consolidados na cultura historiográfica contemporânea.
A cronologia deve ser pensada como um instrumento compartilhado por professores de História com vistas à problematização da proposta, justificação do \textbf{sentido} ( contido no sequenciamento ) e discussão dos \textbf{significados} dos eventos selecionados por diferentes culturas e sociedades.
O ensino de História se justifica na relação do presente com o passado, valorizando o tempo vivido pelo estudante e seu \textbf{protagonismo}, para que ele possa \textbf{participar} ativamente da construção de uma sociedade justa, democrática e inclusiva.

% matched lemmas: forma, participar, protagonismo, sentido, significado
\end{textsample}
